% ==============================================================================
% APPENDIX A: API REFERENCE
% ==============================================================================

\chapter{API Reference}

\section{Authentication Endpoints}

\begin{longtable}{lll}
\caption{Authentication API Endpoints}
\label{tab:auth-api} \\
\toprule
\textbf{Method} & \textbf{Path} & \textbf{Description} \\
\midrule
\endfirsthead
\caption*{Table~\ref{tab:auth-api}: Authentication API Endpoints (continued)} \\
\toprule
\textbf{Method} & \textbf{Path} & \textbf{Description} \\
\midrule
\endhead
\midrule
\endfoot
\bottomline
\endlastfoot
POST & /api/v1/auth/signup & Create a new user account \\
POST & /api/v1/auth/login & Authenticate and receive access token \\
POST & /api/v1/auth/refresh & Refresh current access token \\
POST & /api/v1/auth/logout & Revoke current access token \\
POST & /api/v1/auth/logout-all & Revoke all access tokens \\
GET & /api/v1/auth/tokens & List active access tokens \\
DELETE & /api/v1/auth/tokens/:id & Revoke specific access token \\
POST & /api/v1/auth/password/forgot & Request password reset email \\
POST & /api/v1/auth/password/reset & Reset password with token \\
POST & /api/v1/auth/email/verification-notification & Send verification email \\
POST & /api/v1/auth/email/verify & Verify email with token \\
\end{longtable}

\section{Dashboard Endpoints}

\begin{longtable}{lll}
\caption{Dashboard API Endpoints}
\label{tab:dashboard-api} \\
\toprule
\textbf{Method} & \textbf{Path} & \textbf{Description} \\
\midrule
\endhead
\midrule
\endfoot
\bottomrule
\endlastfoot
GET & /api/v1/dashboard & Unified overview (role-filtered) \\
GET & /api/v1/dashboard/farmer-counts & Farmer counts per organization \\
GET & /api/v1/dashboard/yield-summary & Yield by crop type \\
GET & /api/v1/dashboard/sync-status & Sync queue status per device \\
GET & /api/v1/dashboard/health & API health check \\
\end{longtable}

\section{agricultural product data Endpoints}

\begin{longtable}{lll}
\caption{agricultural product data API Endpoints}
\label{tab:data-api} \\
\toprule
\textbf{Method} & \textbf{Path} & \textbf{Description} \\
\midrule
\endhead
\midrule
\endfoot
\bottomrule
\endlastfoot
GET & /api/v1/farmers & List farmers (paginated) \\
POST & /api/v1/farmers & Create farmer \\
GET & /api/v1/farmers/stats & Farmer statistics \\
GET & /api/v1/farmers/:id & Get farmer details \\
PATCH & /api/v1/farmers/:id & Update farmer \\
DELETE & /api/v1/farmers/:id & Soft-delete farmer \\
POST & /api/v1/farmers/:id/restore & Restore deleted farmer \\
GET & /api/v1/farmers/:id/plots & Get farmer's plots \\
GET & /api/v1/plots & List plots (paginated) \\
POST & /api/v1/plots & Create plot \\
GET & /api/v1/plots/nearby & Find nearby plots \\
GET & /api/v1/plots/:id & Get plot details \\
PATCH & /api/v1/plots/:id & Update plot \\
DELETE & /api/v1/plots/:id & Soft-delete plot \\
POST & /api/v1/plots/:id/restore & Restore deleted plot \\
GET & /api/v1/observations & List observations (paginated) \\
POST & /api/v1/observations & Create observation \\
GET & /api/v1/observations/summary & Observation summary \\
GET & /api/v1/observations/:id & Get observation details \\
PATCH & /api/v1/observations/:id & Update observation \\
DELETE & /api/v1/observations/:id & Soft-delete observation \\
POST & /api/v1/observations/:id/restore & Restore deleted observation \\
PATCH & /api/v1/observations/:id/harvest & Record harvest data \\
\end{longtable}

\section{Machine Learning Endpoints}

\begin{longtable}{lll}
\caption{ML Prediction API Endpoints}
\label{tab:ml-api} \\
\toprule
\textbf{Method} & \textbf{Path} & \textbf{Description} \\
\midrule
\endhead
\midrule
\endfoot
\bottomrule
\endlastfoot
GET & /api/v1/predictions/ml-health & ML service health status \\
POST & /api/v1/predictions & Submit single prediction \\
GET & /api/v1/predictions & List prediction history \\
GET & /api/v1/predictions/:id & Get prediction details \\
GET & /health & (ML service) Service health \\
POST & /predict/yield & (ML service) Single prediction \\
POST & /predict/yield/batch & (ML service) Batch prediction \\
\end{longtable}

\section{Data Import Endpoints}

\begin{longtable}{lll}
\caption{Data Import API Endpoints}
\label{tab:import-api} \\
\toprule
\textbf{Method} & \textbf{Path} & \textbf{Description} \\
\midrule
\endhead
\midrule
\endfoot
\bottomrule
\endlastfoot
POST & /api/v1/imports & Upload CSV/Excel file \\
GET & /api/v1/imports & List import jobs \\
GET & /api/v1/imports/:id & Get import job details \\
GET & /api/v1/imports/:id/schema & Get detected schema \\
GET & /api/v1/imports/:id/records & Get imported records \\
\end{longtable}

\section{RAG Query Endpoints}

\begin{longtable}{lll}
\caption{RAG Query API Endpoints}
\label{tab:rag-api} \\
\toprule
\textbf{Method} & \textbf{Path} & \textbf{Description} \\
\midrule
\endhead
\midrule
\endfoot
\bottomrule
\endlastfoot
GET & /api/v1/rag/status & RAG system health \\
POST & /api/v1/rag/query & Submit natural language query \\
\end{longtable}

\section{Synchronization Endpoints}

\begin{longtable}{lll}
\caption{Sync Queue API Endpoints}
\label{tab:sync-api} \\
\toprule
\textbf{Method} & \textbf{Path} & \textbf{Description} \\
\midrule
\endhead
\midrule
\endfoot
\bottomrule
\endlastfoot
GET & /api/v1/sync-queue & List sync queue entries \\
POST & /api/v1/sync-queue & Create sync entry \\
GET & /api/v1/sync-queue/stats & Sync queue statistics \\
POST & /api/v1/sync-queue/batch & Batch sync submission \\
GET & /api/v1/sync-queue/:id & Get sync entry details \\
PATCH & /api/v1/sync-queue/:id/status & Update sync status \\
POST & /api/v1/sync-queue/:id/retry & Retry failed sync \\
DELETE & /api/v1/sync-queue/:id & Delete sync entry \\
\end{longtable}

\section{Administration Endpoints}

\begin{longtable}{lll}
\caption{Administration API Endpoints}
\label{tab:admin-api} \\
\toprule
\textbf{Method} & \textbf{Path} & \textbf{Description} \\
\midrule
\endhead
\midrule
\endfoot
\bottomrule
\endlastfoot
GET & /api/v1/app-users & List application users \\
POST & /api/v1/app-users & Create application user \\
GET & /api/v1/app-users/:id & Get user details \\
PATCH & /api/v1/app-users/:id & Update user \\
DELETE & /api/v1/app-users/:id & Delete user \\
PATCH & /api/v1/app-users/:id/activate & Activate/deactivate user \\
GET & /api/v1/app-users/:id/devices & List user's devices \\
GET & /api/v1/organizations & List organizations \\
POST & /api/v1/organizations & Create organization \\
GET & /api/v1/organizations/:id & Get organization details \\
PATCH & /api/v1/organizations/:id & Update organization \\
DELETE & /api/v1/organizations/:id & Delete organization \\
GET & /api/v1/organizations/:id/dashboard & Organization dashboard \\
GET & /api/v1/crop-types & List crop types \\
POST & /api/v1/crop-types & Create crop type \\
GET & /api/v1/crop-types/categories & List crop categories \\
GET & /api/v1/crop-types/:id & Get crop type details \\
PATCH & /api/v1/crop-types/:id & Update crop type \\
DELETE & /api/v1/crop-types/:id & Delete crop type \\
GET & /api/v1/devices & List devices \\
POST & /api/v1/devices & Register device \\
GET & /api/v1/devices/by-uuid/:uuid & Find device by UUID \\
GET & /api/v1/devices/:id & Get device details \\
PATCH & /api/v1/devices/:id & Update device \\
PATCH & /api/v1/devices/:id/sync & Mark device synced \\
DELETE & /api/v1/devices/:id & Delete device \\
\end{longtable}

\section{Audit and Monetization Endpoints}

\begin{longtable}{lll}
\caption{Audit and Monetization API Endpoints}
\label{tab:audit-api} \\
\toprule
\textbf{Method} & \textbf{Path} & \textbf{Description} \\
\midrule
\endhead
\midrule
\endfoot
\bottomrule
\endlastfoot
GET & /api/v1/audit-logs & List audit logs \\
POST & /api/v1/audit-logs & Create audit log entry \\
GET & /api/v1/audit-logs/stats & Audit log statistics \\
GET & /api/v1/audit-logs/record/:table/:id & Get audit trail for record \\
GET & /api/v1/audit-logs/:id & Get audit log details \\
GET & /api/v1/monetization/products & List data products \\
POST & /api/v1/monetization/products & Create data product \\
GET & /api/v1/monetization/products/:id & Get product details \\
PATCH & /api/v1/monetization/products/:id & Update product \\
GET & /api/v1/monetization/subscriptions & List subscriptions \\
POST & /api/v1/monetization/subscriptions & Create subscription \\
GET & /api/v1/monetization/transactions & List transactions \\
POST & /api/v1/monetization/chapa/webhook & Chapa payment webhook \\
\end{longtable}

% ==============================================================================
% APPENDIX B: DEPLOYMENT CONFIGURATION
% ==============================================================================

\chapter{Deployment Configuration}

\section{Environment Variables}

The following environment variables are required for system deployment:

\begin{longtable}{ll}
\caption{Required Environment Variables}
\label{tab:env-vars} \\
\toprule
\textbf{Variable} & \textbf{Description} \\
\midrule
\endhead
\midrule
\endfoot
\bottomrule
\endlastfoot
DATABASE\_URL & PostgreSQL connection string \\
ML\_SERVICE\_URL & FastAPI microservice URL \\
ANTHROPIC\_API\_KEY & Claude API key for RAG \\
CEREBRAS\_API\_KEY & Cerebras API key (alternative) \\
OPENROUTER\_API\_KEY & OpenRouter API key (fallback) \\
CHAPA\_SECRET\_KEY & Chapa payment gateway key \\
SESSION\_DRIVER & Session storage driver \\
MAIL\_HOST & SMTP server for email \\
MAIL\_PORT & SMTP port \\
MAIL\_USERNAME & SMTP authentication username \\
MAIL\_PASSWORD & SMTP authentication password \\
REDIS\_HOST & Redis cache host (optional) \\
REDIS\_PORT & Redis cache port (optional) \\
\end{longtable}

\section{Database Migration Commands}

The following commands manage the database schema:

\begin{lstlisting}[caption=Database Migration Commands,language=bash]
# Run all pending migrations
node ace migration:run

# Rollback last migration batch
node ace migration:rollback

# Run all migrations from scratch
node ace migration:run --force

# Create a new migration file
node ace make:migration table_name

# Refresh database views
node ace migration:refresh
\end{lstlisting}

\section{ML Service Training Commands}

The following commands train and manage ML models:

\begin{lstlisting}[caption=ML Service Training Commands,language=bash]
# Train using PostgreSQL database (default)
python -m ml_service.train --source db

# Train using CSV file
python -m ml_service.train --source csv \
  --data ml-service/data/sample_yield.csv

# Train with custom hyperparameters
python -m ml_service.train \
  --n-estimators 300 \
  --max-depth 7 \
  --lr 0.08

# Start ML microservice for inference
python -m ml_service

# Run ML service tests
python -m pytest ml-service/tests/
\end{lstlisting}

% ==============================================================================
% APPENDIX C: SAMPLE TRAINING DATA
% ==============================================================================

\chapter{Sample Training Data}

\section{Training Data Format}

The training data is stored in CSV format with the following columns:

\begin{longtable}{ll}
\caption{Training Data Column Descriptions}
\label{tab:train-cols} \\
\toprule
\textbf{Column} & \textbf{Description} \\
\midrule
\endhead
\midrule
\endfoot
\bottomrule
\endlastfoot
crop\_name & Crop type (teff, wheat, maize, barley, sorghum, sesame, chickpea, coffee) \\
season & Growing season (meher, belg) \\
region & Administrative region (Afar, Amhara, Oromia, Sidama, SNNPR, Tigray) \\
year & Harvest year (2015--2024) \\
area\_hectares & Cultivated area in hectares \\
rainfall\_mm & Annual rainfall in millimeters \\
temperature\_celsius & Average growing season temperature \\
fertilizer\_amount\_kg & Fertilizer applied in kilograms \\
actual\_yield\_kg & Actual harvested yield in kilograms (target) \\
\end{longtable}

\section{Sample Records}

The first ten records from the training dataset illustrate the data format and range of values:

\begin{table}[H]
\centering
\caption{Sample Training Records}
\label{tab:sample-records}
\begin{tabular}{llllrrrrr}
\toprule
\textbf{Crop} & \textbf{Season} & \textbf{Region} & \textbf{Year} & \textbf{Area} & \textbf{Rain} & \textbf{Temp} & \textbf{Fert} & \textbf{Yield} \\
\midrule
maize & meher & Afar & 2019 & 5.28 & 509.3 & 12.6 & 148.1 & 14242.5 \\
teff & meher & Oromia & 2018 & 5.04 & 1203.0 & 24.0 & 143.2 & 9599.5 \\
sesame & meher & Tigray & 2024 & 5.92 & 1604.0 & 29.0 & 31.9 & 5454.9 \\
sorghum & belg & Oromia & 2016 & 7.91 & 838.5 & 18.6 & 52.9 & 17479.3 \\
teff & belg & Sidama & 2016 & 19.48 & 867.8 & 23.8 & 165.9 & 30890.9 \\
sorghum & meher & Oromia & 2018 & 15.57 & 1777.8 & 31.4 & 173.3 & 35483.0 \\
sesame & belg & Tigray & 2020 & 3.67 & 832.9 & 26.8 & 140.4 & 3581.1 \\
wheat & meher & Amhara & 2022 & 7.90 & 1784.3 & 26.0 & 111.4 & 22795.9 \\
coffee & meher & Amhara & 2015 & 16.20 & 901.7 & 11.7 & 182.6 & 16957.4 \\
sorghum & belg & Tigray & 2022 & 3.29 & 509.4 & 28.6 & 107.8 & 6115.2 \\
\bottomrule
\end{tabular}
\end{table}

The dataset contains 2,000 records with the following distribution characteristics:
\begin{itemize}
  \item Area ranges from approximately 3 to 20 hectares
  \item Rainfall ranges from approximately 500 to 1800 mm
  \item Temperature ranges from approximately 11 to 33 degrees Celsius
  \item Fertilizer ranges from approximately 30 to 185 kg
  \item Yield ranges from approximately 3,500 to 85,000 kg
\end{itemize}

\section{Data Quality Notes}

The training data was cleaned and preprocessed before model training:
\begin{itemize}
  \item Records with missing target (actual\_yield\_kg) values were excluded
  \item Infinite values were replaced with NaN and imputed using column median
  \item Numerical features were normalized using z-score standardization
  \item Categorical features were encoded using label encoding
  \item Season names were normalized to lowercase
  \item The final dataset passed all validation checks with a 100 percent valid row rate
\end{itemize}
